\section{Filling punctures and reconstructing the complete curve}\label{part:C}

The three rows in \S32, folio 202, scan 399 must be read separately.
They concern an arithmetic image, a geometric completion, and the
full group of loop-preserving outer automorphisms.  The first two
have exact filling sequences.  The unrestricted third row, marked
``??'', requires a further assertion about the full automorphism
group.  It is not a formal consequence of either of the first two.
Throughout this part a filling map from type $(g,n+1)$ has hyperbolic
target $(g,n)$ unless the exceptional case is explicitly discussed.

\subsection{The quotient and the two geometric exact sequences}

\statement{C.1}{Proved}{Filling is quotient by the closed normal inertia subgroup}
Let $U=X\setminus S$ be a smooth geometrically connected curve over
a field $K$ of characteristic zero, with smooth proper completion
$X$, and let $S_0\subset S$ be a Galois-stable subset.  Put
$U'=X\setminus(S\setminus S_0)$.  After choosing compatible
geometric base points, the inclusion $U\hookrightarrow U'$ induces
a continuous epimorphism
\[
 \pi_1^{\mathrm{\acute et}}(U)\longrightarrow
 \pi_1^{\mathrm{\acute et}}(U')
\]
whose kernel is the closed normal subgroup generated by the
geometric inertia groups at the points of $S_0$.  The same statement
holds for the geometric groups, and the quotient is compatible
with the maps to $G_K$.  In particular, the quotient classifies
precisely the finite \emph{\'etale} coverings extending across the
filled punctures.

\begin{proof}
Normalize $U'$ in the function field of a connected finite \'etale
cover of $U$.  Normalization is finite for curves of finite type
over a field.  The resulting finite morphism is \'etale at a point
of $S_0$ exactly when its inertia action on the associated finite
set is trivial; in characteristic zero no wild-inertia qualification
is needed.  Thus finite \'etale covers of $U'$ identify with the
full subcategory of finite continuous $\pi_1(U)$-sets on which all
these inertia groups act trivially.  The fundamental group of
this Galois category is the quotient by their closed normal
closure.  The same argument after base change to $\overline K$
gives the geometric assertion.  The structural maps to
$\operatorname{Spec}K$ commute with the inclusion, giving
compatibility with $G_K$.  A change of geometric base point or
connecting path changes the homomorphisms by the corresponding
inner automorphisms of the target.  This is the group-theoretic
filling operation used in \S28, folios 175--176, scans 347--348,
and \S32, folios 200--202, scans 395--400.
\end{proof}

\statement{C.2}{Proved}{The full profinite Birman sequence}
For $2g-2+n>0$, forgetting the last marked point gives a short exact
sequence of profinite groups
\[
 1\longrightarrow\pi_{g,n}
 \longrightarrow\widehat P_{g,n+1}
 \longrightarrow\widehat P_{g,n}\longrightarrow1.
 \tag{C.1}
\]
The first arrow is the completed point-pushing homomorphism.
Allowing permutations of the retained punctures replaces the two
mapping class groups by $M^+_{g,n+1,i}$ and $M^+_{g,n}$, where the
source fixes the puncture being forgotten.  The same version holds
with orientation reversal allowed.

\begin{proof}
The discrete Birman sequence is
$1\to\Pi_{g,n}\to P_{g,n+1}\to P_{g,n}\to1$.
We justify the potentially delicate passage to full completions.
The action by conjugation on its normal subgroup $\Pi_{g,n}$
gives $P_{g,n+1}\to\operatorname{Aut}(\pi_{g,n})$, hence a map
from its profinite completion.  The composite
\[
 \pi_{g,n}\longrightarrow\widehat P_{g,n+1}
 \longrightarrow\operatorname{Aut}(\pi_{g,n})
\]
is the inner-automorphism map, which is injective because
$\pi_{g,n}$ is center-free.  Thus the first arrow of (C.1) is
injective.  Completion is right exact in the following precise
sense: $\widehat P_{g,n+1}$ modulo the closure of the point-pushing
subgroup is $\widehat P_{g,n}$, because both classify the same
finite quotients trivial on that subgroup.  The image of
$\pi_{g,n}$ is compact and contains the dense discrete
point-pushing group, so it is exactly that closure.  This proves
exactness.  The finite permutation and orientation quotients give
the other versions by the identical argument.  The exact sequence
is also explicitly recorded in
\citet[\S2, p.~4]{Boggi2009BC}.
\end{proof}

The hyperbolic target hypothesis is essential.  Filling the sole
puncture of a torus induces the isomorphism
$P_{1,1}\simeq P_{1,0}\simeq\operatorname{SL}_2(\mathbb Z)$;
the alleged kernel $\pi_{1,0}=\widehat{\mathbb Z}^2$ would be
wrong.  Point-pushing is killed by translations isotopic to the
identity on the unmarked torus.  Similarly,
$P_{0,3}\to P_{0,2}$ is a map of trivial groups, not an extension
with kernel $\pi_{0,2}=\widehat{\mathbb Z}$.  For targets $(0,0)$
and $(0,1)$ the pure oriented groups and geometric fundamental
groups are trivial, so the corresponding trivial sequence carries
no anabelian content.  For hyperbolic types the \emph{surface}
group kernel is center-free, even in cases where a \emph{mapping
class} group has center: the elliptic involution at $(1,1)$ and the
hyperelliptic involution at $(2,0)$ are retained.  For the pure
procongruence groups the exact center statement is
\citet[Theorem 4.4]{Boggi2020BC}: the center equals the discrete
center, of order two in these two cases and trivial in the other
hyperbolic pure cases.  Additional centers in variants permitting
permutations must not be suppressed by this pure-group statement.

\statement{C.3}{Known}{The procongruence Birman sequence}
For every hyperbolic target type there is an exact sequence
\[
 1\longrightarrow\pi_{g,n}
 \longrightarrow\check P_{g,n+1}
 \longrightarrow\check P_{g,n}\longrightarrow1.
 \tag{C.2}
\]
It is compatible with (C.1), and the map on the left-hand kernels
is the identity under their natural identifications.  The
point-pushing subgroup therefore has its full profinite topology
inside the procongruence completion.

This is \citet[Corollary 4.7, with $\mathscr C=\mathscr F$,
and Proposition 2.4(ii)]{Boggi2020BC}.  Its Remark 4.8 specifies
the history: Matsumoto's theorem treated open targets, and
Hoshi--Mochizuki extended it to closed targets.  The earlier
attached \citet[Theorem 2.5]{Boggi2009BC} alone has the extra
condition $n\geq1$ and cannot be cited for filling the last
puncture of a higher-genus surface.  Corollary 4.7 also has a
pro-$p$ congruence form, with kernel the pro-$p$ surface group,
for every prime $p$, including $2$ and $3$.  This does not identify
an arbitrary pro-$p$ congruence completion with the pro-$p$
completion of the whole discrete mapping class group.

\statement{C.4}{Proved}{The congruence kernel is unchanged by puncturing}
The natural map
\[
 C_{g,n+1}\longrightarrow C_{g,n}
\]
is an isomorphism for every hyperbolic target $(g,n)$.

\begin{proof}
Compare (C.1) and (C.2).  If an element of $C_{g,n+1}$ maps to
the identity of $\widehat P_{g,n}$, it lies in $\pi_{g,n}$; this
subgroup maps injectively to the lower middle group, so the
element is trivial.  Conversely, take $a\in C_{g,n}$ and lift it
to $b\in\widehat P_{g,n+1}$.  Its image in
$\check P_{g,n+1}$ lies in the left-hand copy of $\pi_{g,n}$.
Multiply $b$ by the inverse of the corresponding element of the
upper copy.  The resulting lift lies in $C_{g,n+1}$.  The map is
therefore a continuous bijection between profinite groups and
hence a topological isomorphism.  This proves exactly which
congruence property transfers in either direction under filling;
it does not assume its truth.
\end{proof}

\subsection{The arithmetic image and the revised quotient}

Let $\mathcal M_{g,n}$ now denote the stack of smooth curves with
\emph{ordered} markings over $\mathbb Q$.  With compatible
geometric base points put
\[
 \mathcal A_{g,n}=\pi_1^{\mathrm{\acute et}}(\mathcal M_{g,n}),
 \quad
 \rho_{g,n}:\mathcal A_{g,n}\to O_{g,n},
 \quad
 N^{\mathrm{arith},!}_{g,n}=\operatorname{im}(\rho_{g,n}).
\]
Its geometric image is $\check P_{g,n}$.  The image is closed
because its source is profinite.  This is the ordered version of
the revised definition in \S29, folio 188, scan 371, not the
normalizer definition of \S26.

\statement{C.5}{Proved}{The revised arithmetic quotient is the absolute Galois group}
For every hyperbolic type, the structural quotient induces a
canonical isomorphism of profinite groups
\[
 G_{\mathbb Q}\ \xrightarrow{\ \sim\ }\
 N^{\mathrm{arith},!}_{g,n}/\check P_{g,n}.
 \tag{C.3}
\]
Moreover, the kernel of the universal monodromy representation
$\rho_{g,n}$ is exactly $C_{g,n}$, viewed in its geometric
subgroup $\widehat P_{g,n}$.  With unordered markings the
corresponding assertion is
$N^{\mathrm{arith}}_{g,n}/\Sigma_{g,n}\simeq G_{\mathbb Q}$.

\begin{proof}
The standard geometric/arithmetic exact sequence of the moduli
stack gives
\[
 1\to\widehat P_{g,n}\to\mathcal A_{g,n}
 \to G_{\mathbb Q}\to1.
\]
Consequently its quotient map to the right-hand side of (C.3)
is surjective by definition.  Let its kernel be $J$.
Choose a marked curve of this type over a finite extension
$L/\mathbb Q$.  Such a curve exists by choosing a geometric
curve and its finitely many markings over $\overline{\mathbb Q}$
and descending their finite defining data; no assertion about
arbitrarily many rational points over $\mathbb Q$ is needed.
The associated point of the moduli stack supplies a section over
$G_L$.  \citet[Corollary 6.4, pp.~66--67]{HoshiMochizuki2011}
says that the preimage in $G_L$ of the geometric monodromy image
is trivial.  Hence $J\cap G_L=1$, so $J$ is a finite normal
subgroup of $G_{\mathbb Q}$.

There is no nontrivial finite normal subgroup of $G_{\mathbb Q}$.
Indeed, a finite normal subgroup is centralized by an open
subgroup, whereas every open subgroup of $G_{\mathbb Q}$ is
slim and has trivial centralizer in $G_{\mathbb Q}$; see
\citet[Corollary 2.2 in the published version; Corollary 1.3
in the author preprint]{Mochizuki2006} and D.14.
Thus $J=1$.  The kernel of $\rho_{g,n}$ is now contained in
the geometric subgroup, where it is by definition $C_{g,n}$.
The unordered version follows by the finite permutation
extension; alternatively its pure subgroup and geometric image
give the same quotient, precisely as in \S32, folio 200,
scan 395, equations (1)--(3).
\end{proof}

The required arithmetic input is a published anabelian theorem.
No section conjecture and no conjectural identification of
$G_{\mathbb Q}$ with $\widehat{\mathrm{GT}}$ is used.
Corollary 6.5 of the same Hoshi--Mochizuki paper states the useful
consequence that total universal monodromy is faithful if and only
if its geometric restriction is faithful.  Thus C.5 does not
silently solve the high-genus congruence problem.

\statement{C.6}{Proved}{The arithmetic filling sequence and invariance of the quotient}
For every hyperbolic target there is a continuous short exact
sequence
\[
 1\longrightarrow\pi_{g,n}
 \longrightarrow N^{\mathrm{arith},!}_{g,n+1}
 \longrightarrow N^{\mathrm{arith},!}_{g,n}
 \longrightarrow1.
 \tag{C.4}
\]
The quotient map on the groups in (C.3) is the identity of
$G_{\mathbb Q}$.  The corresponding sequence with retained
markings unordered uses the stabilizer of the forgotten marking
in the source.  In particular, the faithfulness condition (f) and
the filling isomorphism of the revised arithmetic quotients in
\S32, folios 201--202, scans 397--400 hold for every hyperbolic
type.

\begin{proof}
The universal punctured curve over $\mathcal M_{g,n}$ is
$\mathcal M_{g,n+1}$.  The fundamental-group homomorphism fits
into the arithmetic exact sequence
\[
 1\to\pi_{g,n}\to\mathcal A_{g,n+1}
 \to\mathcal A_{g,n}\to1,
\]
whose geometric row is (C.1) and whose quotient row is the
identity of $G_{\mathbb Q}$.  One may obtain exactness directly
by this diagram: lift an element over a prescribed Galois
element, then adjust by geometric surjectivity; the kernel lies
in the geometric row.  Universal monodromy is compatible with
filling, so quotient this sequence by its monodromy kernels.
By C.5 they are $C_{g,n+1}$ and $C_{g,n}$, and by C.4 the map
between them is an isomorphism.  Their intersection with
$\pi_{g,n}$ is trivial by C.3.  It follows that the quotient
sequence is precisely (C.4).  The arithmetic quotients in the
resulting diagram are induced by the same structural morphism
to $\operatorname{Spec}\mathbb Q$, hence their comparison is
canonically the identity on $G_{\mathbb Q}$.

The same argument over a finite extension $L$ gives the identity
on $G_L$, compatibly with restriction from $G_{\mathbb Q}$.
Changing the base points conjugates the group homomorphisms by
geometric elements.  It has no effect on their arithmetic
quotients, which is the precise sense in which their isomorphism
is canonical.  For coverings defined over a number field the
arithmetic quotients likewise compare through that same field
group.  This last observation concerns arithmetic images; it
does not identify the full normalizers of arbitrary geometric
structures under every abstract finite covering.
\end{proof}

On the scan of page 399, the first sequence has $N$, rather than
the based automorphism group $M$, and its arrow is directed
towards the type with one fewer puncture.  Equation (C.4) is its correct
arithmetic interpretation.  On scan 395, equation (4) omits the
pure-marking superscript on the geometric denominator in its
left-hand quotient; equation (3) immediately above supplies the
mathematically consistent version.  We record both observations
as evidence for the editors, not as silent changes to the source.

\subsection{What the sequence marked ``??'' requires}

\statement{C.7}{Proved}{The canonical based map supplied by a distinguished cusp}
Fix a representative inertia subgroup $I_i\subset\pi_{g,n+1}$
and put $\pi'=\pi_{g,n+1}/\overline{\langle I_i\rangle}^{\mathrm{normal}}$
of type $(g,n)$.  There is a continuous homomorphism
\[
 \theta_i:O_{g,n+1,i}\longrightarrow
 \operatorname{Aut}_{\mathrm{lac}}(\pi'),
 \tag{C.5}
\]
where $O_{g,n+1,i}$ stabilizes the conjugacy class of $I_i$.
Its composite with the outer quotient is the usual filling map
$f_i:O_{g,n+1,i}\to O_{g,n}$.

\begin{proof}
Represent an outer automorphism by an actual automorphism
preserving $I_i$, which is possible by the definition of its
stabilizer.  It induces an automorphism of the quotient $\pi'$.
Two such representatives differ by conjugation by an element
of $N_{\pi_{g,n+1}}(I_i)=I_i$; the equality follows from
\citet[Proposition 3.6]{BoggiZalesskii2019BC}.  That conjugation
is killed in $\pi'$.  Thus the induced genuine automorphism is
independent of the representative.  Composition and continuity
follow by passing from the closed stabilizer of $I_i$ in the
automorphism group to its compact quotient.  Replacing $I_i$ by
a conjugate conjugates the construction by the corresponding
inner automorphism of $\pi'$.  This is the based ambiguity.
On the point-pushing subgroup
$\pi_{g,n}\subset\check P_{g,n+1}\subset O_{g,n+1,i}$,
the map is the inner-automorphism embedding of $\pi_{g,n}$ in
$\operatorname{Aut}(\pi')$, as is seen first in the discrete
Birman sequence and then by continuity.
\end{proof}

\statement{C.8}{Open}{The unrestricted filling assertion and its exact equivalent forms}
With the choices of C.7, the unrestricted third row of scan 399
would assert
\[
 1\longrightarrow\pi_{g,n}
 \longrightarrow O_{g,n+1,i}
 \xrightarrow{f_i}O_{g,n}\longrightarrow1.
 \tag{C.6?}
\]
This is equivalent to $\theta_i$ in (C.5) being an isomorphism.
It is also equivalent to bijectivity of the induced filling map
on the unrestricted oriented prediscretization spaces
\[
 O_{g,n+1,i}/\Sigma_{g,n+1,i}
 \longrightarrow O_{g,n}/\Sigma_{g,n}.
 \tag{C.7?}
\]
Here the source is canonically the corresponding marking
quotient with the chosen cusp filled: the geometric permutation
group permits the cusp to be fixed before taking the quotient.
For individually labeled punctures use the pure variants in
both numerator and denominator.  We have not proved (C.6?)
for the full $O$ groups.  Its smallest hyperbolic target instance
is $(0,4)\to(0,3)$.  The remaining conditions in this instance,
and in general, are explicitly
\[
 f_i(O_{g,n+1,i})=O_{g,n},\qquad
 \ker(f_i)=\pi_{g,n}.
\]

\begin{proof}[Proof of the stated equivalences]
Since $\pi'$ is center-free, the lower row of
\[
\begin{array}{ccccccccc}
1&\to&\pi_{g,n}&\to&O_{g,n+1,i}&\xrightarrow{f_i}&O_{g,n}&\to&1\\
 &&\Vert&&\downarrow\theta_i&&\Vert\\
1&\to&\operatorname{Inn}(\pi')&\to&
\operatorname{Aut}_{\mathrm{lac}}(\pi')&\to&O_{g,n}&\to&1
\end{array}
\]
is exact, with the left identification supplied by C.7.
Exactness of the upper row is consequently equivalent to
$\theta_i$ being an isomorphism, by elementary lifting and
kernel arguments.  C.3 with permutations gives
$f_i(\Sigma_{g,n+1,i})=\Sigma_{g,n}$ and kernel
$\pi_{g,n}$ on that geometric subgroup.  By B.13, (C.7?) is
bijective if and only if $f_i$ is onto and
$f_i^{-1}(\Sigma_{g,n})=\Sigma_{g,n+1,i}$.  Under geometric
surjectivity, the latter is equivalent to
$\ker(f_i)\subset\Sigma_{g,n+1,i}$, and then the known
geometric kernel makes it equivalent to the second condition
above.  This proves all the equivalences without assuming any
of them.
\end{proof}

Thus the two geometric rows and the arithmetic row have been
settled, while the unrestricted double-hat row has been reduced
to its exact missing content.  The map to genuine automorphisms
in scan 348 and the map to outer automorphisms in scan 399 are
not interchangeable: the former would be an isomorphism; the
latter would have the point-pushing kernel.  The manuscript's
``peut-\^etre enti\`erement faux'' in folio 176, scan 349 is a
question, not a counterexample.  As in B.7 and B.14, the label
Open here specifies what this audit has not established; no
claim is made that a comprehensive literature proof of
independence has been supplied.  In particular, we do not
rename this unresolved lifting problem the congruence subgroup
problem: both involved genus-zero congruence assertions are known.

\statement{C.9}{Known}{The established combinatorial cuspidalization theorem}
Let $X$ be a hyperbolic curve over an algebraically closed field
of characteristic zero, and let $X_m$ be its configuration space
of $m$ ordered distinct points.  Let $\Lambda$ be either all
primes or a singleton, and let $\pi_m$ be the maximal pro-$\Lambda$
quotient of $\pi_1(X_m)$.  Define
$\operatorname{Out}^{\mathrm{FC}}(\pi_m)$ by preservation of
the fiber subgroups of all configuration projections and of the
cuspidal inertia classes in their surface-group subquotients,
as in the cited definition.  Then forgetting one point gives
\[
 \operatorname{Out}^{\mathrm{FC}}(\pi_{m+1})
 \longrightarrow\operatorname{Out}^{\mathrm{FC}}(\pi_m)
\]
which is injective for all $m\geq1$, bijective for $m\geq3$ if
$X$ is affine, and bijective for $m\geq4$ if $X$ is proper.
The factor-permutation subgroup centralizes these FC-admissible
outer automorphisms.

These are the exact assertions of
\citet[Theorem 6.1, pp.~61--63]{HoshiMochizuki2011}, extending
\citet[Theorem 4.1]{Mochizuki2010BC}.  In the original definition,
F-admissibility preserves each fiber subgroup, rather than
merely the collection up to an unspecified permutation;
FC-admissibility adds cuspidal compatibility on the relevant
subquotients.  The image of an F-admissible automorphism from
one higher configuration level is FC-admissible, by
\citet[Theorem A and Corollary 1.9]{HoshiMochizukiTopicsBC}.
These extra data are a substantive hypothesis.  A single
filled surface group is a fiber in this construction, not the
whole configuration-space group, and an arbitrary member of
$O_{g,n+1,i}$ has not thereby been supplied with such a lift.
Consequently this theorem does not by itself remove the
question mark in (C.6?).  Its full and single-prime forms include
the primes $2$ and $3$ in characteristic zero; its positive
characteristic version requires the characteristic to lie
outside the chosen set of primes.

\subsection{What it means to reconstruct the entire \texorpdfstring{\'etale}{etale} topos}

\statement{C.10}{Refuted}{The finite-cover topos is not the full \texorpdfstring{\'etale}{etale} topos}
For a connected positive-dimensional smooth curve $X$, the
topos of continuous discrete $\pi_1^{\mathrm{\acute et}}(X)$-sets
need not, and in fact cannot, be its full small \'etale topos.
Thus C.1 alone does not reconstruct ``tout le topos \'etale''.

\begin{proof}
In the classifying topos of any group, a subobject of the
terminal object is an invariant subset of a one-element set.
There are exactly two such subobjects.  In the small \'etale
topos of a curve, every Zariski open subset represents a
subterminal object.  A nonempty proper open subset gives a
third one, distinct from the empty and the whole scheme.
An equivalence of topoi preserves subterminal objects, so the
topoi cannot be equivalent.  One particularly transparent case
is $X=\mathbb P^1_{\overline{\mathbb Q}}$: its geometric
fundamental group is trivial, while its \'etale topos still
contains all its distinct open subobjects.  The finite
two-versus-three subobject count is also recorded in
\texttt{code/bc\_counterexamples.py}.  The argument applies to
the classifying topos of all continuous discrete sets, so it
does not rely on confusing finite sets with all such sets.
\end{proof}

\statement{C.11}{Proved}{Reconstruction of the complete curve and its full \texorpdfstring{\'etale}{etale} topos}
Fix a sub-$p$-adic field $K$ of characteristic zero, in
particular a number field or a finite extension of
$\mathbb Q_p$.  On the essential image of hyperbolic curves
over $K$, the arithmetic extension
\[
 1\longrightarrow\pi_1(U_{\overline K})
 \longrightarrow\Pi_U=\pi_1(U)
 \longrightarrow G_K\longrightarrow1
 \tag{C.8}
\]
reconstructs the smooth proper completion $X$ of $U$ and its
entire small \'etale topos, functorially for $K$-isomorphisms.
The group datum includes the projection to the specified
$G_K$; isomorphisms of it are taken modulo inner automorphisms
from the geometric kernel.  The resulting topos is determined
up to the equivalence induced by the reconstructed
$K$-isomorphism, with the usual natural-isomorphism ambiguity.

\begin{proof}
The relative anabelian Isom theorem in D.4 follows from
\citet[Theorem A, equivalently Theorem 16.5 and its following
Remark]{Mochizuki1999}: for hyperbolic curves over a sub-$p$-adic
field, the map from dominant $K$-morphisms to continuous open
homomorphisms of arithmetic fundamental groups over $G_K$,
modulo geometric inner automorphisms of the target, is
bijective.  Applied to an isomorphism and its inverse, this
reconstructs a unique $K$-isomorphism $U\simeq U'$ inducing
the prescribed outer class.  Consequently the functor on
isomorphism groupoids is fully faithful, and an inverse on
its essential image reconstructs $U$ with precisely this
uniqueness.  Existence is a hypothesis on the input datum,
not a claim that every extension of $G_K$ is geometric.

A smooth geometrically connected curve has a unique smooth
proper completion up to unique isomorphism extending the
identity on $U$.  It may be constructed as the nonsingular
proper model of its one-variable function field $F=K(U)$.
An isomorphism of open curves extends uniquely to these
proper models.  It remains to show explicitly that this
reconstructs the whole topos rather than just its finite
coverings.

For every finite separable field extension $L/F$, let $X_L$
be the normalization of $X$ in $L$.  It is finite over $X$.
Consider the site whose objects are finite disjoint unions
of opens $V\subset X_L$ on which $V\to X$ is \'etale.
Morphisms are $X$-morphisms, equivalently the compatible
generic field embeddings and their extensions on the chosen
opens.  Covering families are jointly surjective \'etale
families.  All these data are determined by the recovered
curve $X$, and hence by (C.8).

This site is a basis for the small \'etale site of $X$.
Indeed, a connected separated \'etale scheme $Y$ of finite
presentation over $X$ is normal and integral, and its
function field is a finite separable extension $L/F$.
Zariski's main theorem identifies $Y$ with an open subscheme
of $X_L$.  One can use the precise factorization of
\citet[Lemmas 37.43.2 and 37.43.3, Tags 02LR and 05K0]{StacksBC}:
a separated quasi-finite morphism factors as an open
immersion followed by a finite morphism; normalizing and
using the same generic field gives the specified
normalization.  Every \'etale scheme over $X$ is covered
by separated affine \'etale opens, and the quasi-compact
ones are of finite presentation.  Hence these objects
form a covering basis even when a general object of the
small site is not separated or quasi-compact.  The
comparison lemma for sites now identifies its sheaf
topos with $\operatorname{Sh}(X_{\mathrm{\acute et}})$.
Transport of normalized extensions and opens along
$K$-isomorphisms proves functoriality.
\end{proof}

This theorem gives a complete affirmative formulation of the
reconstruction sought in \S32, folios 206--217, scans 407--430.
The exact phrase quoted in the prompt was not located as a
literal sentence in scans 395--420; the subsequent construction
of varying open groupoids and coverings occurs in scans
421--430.  We cite that development rather than manufacture
a quotation.  The indispensable input in C.11 is the
\emph{arithmetic} group over its field, and the proof uses
relative anabelian reconstruction before constructing the
site.  A bare geometric $\pi_{g,n}$ is the same for every
complex curve of that type and cannot recover its algebraic
moduli.  Nor does retaining only the finite covers unramified
over the original $U$ produce the displayed site: finite
extensions branched at newly removed points are also needed.
We do not claim that the manuscript's direct construction
from an independently characterized class of admissible
splittings has thereby been completed without anabelian
reconstruction.

\statement{C.12}{Proved}{The correct point-filling statement over a field}
For a smooth curve $U'$ over a field $K$, a $K$-rational point
$x\in U'(K)$, and $U=U'\setminus\{x\}$,
\[
 U(K)=U'(K)\setminus\{x\}.
\]
Under the anabelian point maps this gives the corresponding
bijection of \emph{geometric point-induced} conjugacy classes
of splittings, compatibly with filling.  If the deleted set
is merely Galois-stable, the correct statement uses its
$K$-rational points; a point of that set need not itself be
rational.  The statement does not assert that every group
section is point-induced.

\begin{proof}
A morphism $\operatorname{Spec}K\to U'$ factors through the
open complement precisely when its image is not $x$.
The equality follows.  The maps on arithmetic fundamental
groups commute with these morphisms, so the same equality
holds for their images in the set of section classes;
injectivity of the point map in the stated anabelian setting
is the theorem recorded in D.5.  Nothing in this argument
requires surjectivity of the section map.  This supplies
the field-valued form of the expectation in \S32, folios
211--212, scans 417--420, while keeping the word
``admissible'' tied to actual geometric points.
\end{proof}

\statement{C.13}{Refuted}{Deleting one section is not deleting one element of a section set}
The analogue of C.12 for sections over an arbitrary base is
false: distinct sections of a family can intersect.  Thus
the assertion in scan 418 cannot be imposed on arbitrary
profinite acting groups by treating all sections as a
set of disjoint geometric points.

\begin{proof}
Over $\mathbb Q$ put
$S=\mathbb P^1\setminus\{0,1,\infty\}$ and consider the
constant hyperbolic family $S\times S\to S$, projection to
the second factor.  The sections
\[
 s(t)=(t,t),\qquad e(t)=(2,t)
\]
are distinct and intersect at $t=2$.  Delete the image of
$e$ from the family.  Its fibers now have type $(0,4)$,
but $s$ does not lift to a section of the deleted family:
at $2$ its value is exactly the deleted point.  It is
therefore false that the new section set is the old one
minus the single section $e$.  All objects and morphisms
are defined over $\mathbb Q$, and the intersection is
checked by substituting $t=2$ in $t-2$ in the accompanying
code.  This also gives a geometric example for the acting
group $\pi_1(S)$, so the defect is not an artifact of
inventing unrelated group actions.

The manuscript itself recognizes this obstruction later:
\S32, folios 219--220, scans 433--435 withdraws the
overgeneralized formulation because sections of families
can meet.  We preserve that withdrawal.  C.12 over a
field and C.11 via the reconstructed curve are the
valid replacements needed for the present task.
\end{proof}
